% =====================================================================
%  CHAPTER 6: 营养与疾病 — 参考：0613笔记 (PRIMARY SOURCE)
% =====================================================================
\chapter{营养与疾病}
\chapterintro{
\keyword{掌握}超重肥胖饮食防治策略原则和方法；掌握心血管疾病、糖尿病、癌症和痛风的营养防治原则。\keyword{熟悉}体重控制意义及方法；熟悉体力活动指南；熟悉临床营养筛查与营养评价基本方法；熟悉病人膳食分类；熟悉肠内营养和肠外营养概念、方法、适应证和并发症；熟悉FSMP概念。\keyword{了解}膳食营养与慢性病的关系；了解临床营养治疗目的和主要方法；了解肠内营养制剂、肠外营养制剂及FSMP分类。
}

% =====================================================================
%  膳食营养与慢性病的关系
% =====================================================================
\section{膳食营养与慢性病的关系}

\begin{tipbox}
\textbf{膳食营养与慢性非传染性疾病（NCD）的关系（了解）：}

慢性非传染性疾病（主要包括心血管疾病、糖尿病、癌症和慢性呼吸系统疾病）是全球最主要的死亡原因。不合理膳食是NCD最重要的可改变危险因素之一。

\textbf{膳食因素与主要慢性病的关联：}
\begin{enumerate}[leftmargin=*]
    \item \textbf{超重/肥胖：}高能量密度食物、含糖饮料、大份量和久坐 → 能量正平衡 → 肥胖 → 是心血管疾病、糖尿病、多种癌症的共同危险因素。
    \item \textbf{心血管疾病：}高钠（高血压）、高饱和脂肪和反式脂肪（动脉粥样硬化）、低蔬果和全谷物（膳食纤维和抗氧化物质不足）。
    \item \textbf{2型糖尿病：}高GI/GL膳食、含糖饮料、红肉和加工肉类摄入过多，全谷物和膳食纤维摄入不足。
    \item \textbf{癌症：}加工肉类（IARC Ⅰ类致癌物）、红肉（ⅡA类）、酒精（多种癌症）、黄曲霉毒素污染食物；保护因素：蔬果、全谷物和膳食纤维。
\end{enumerate}

\textbf{公共健康膳食策略的核心理念：}
合理膳食是NCD一级预防的基石，膳食因素可干预、可改变，通过人群层面的营养政策和个体层面的营养教育可实现NCD的有效预防。膳食营养与慢性病的关系为制定膳食指南和营养政策提供了重要的科学依据。
\end{tipbox}

% =====================================================================
%  营养与肥胖
% =====================================================================
\section{营养与肥胖}

\begin{defbox}
\textbf{肥胖的定义与判定标准：}
\begin{itemize}[leftmargin=*]
    \item \textbf{肥胖：}体内脂肪过多积累，能量摄入长期大于能量消耗所致。
    \item \textbf{体质指数（BMI）：}\imp{BMI = 体重(kg) / [身高(m)]$^2$}。我国标准：BMI $<$ 18.5 消瘦，18.5$\sim$23.9 正常，24.0$\sim$27.9 超重，\imp{BMI $\geq$ 28 肥胖}。
    \item \textbf{中心性肥胖：}腰围\imp{男 $\geq$ 85cm，女 $\geq$ 80cm}，或腰臀比男 $>$ 0.9、女 $>$ 0.85。
\end{itemize}
\end{defbox}

\begin{keybox}
\textbf{肥胖的膳食因素与预防控制：}

\textbf{膳食因素：}
\begin{itemize}[leftmargin=*]
    \item 高能量密度食物（油炸食品、甜点、快餐）
    \item 含糖饮料（添加糖占总能量 $>$ 10\%）
    \item 进食量大、进餐速度快、不吃早餐、夜间进食
    \item 膳食纤维摄入不足，蔬果和全谷物摄入不足
\end{itemize}

\textbf{能量平衡原理：}肥胖的本质是\imp{能量摄入长期大于能量消耗}。能量消耗包括基础代谢（约60\%$\sim$70\%）、身体活动（约15\%$\sim$30\%）和食物热效应（约10\%）。

\textbf{预防控制措施：}
\begin{enumerate}[leftmargin=*]
    \item \textbf{控制总能量摄入：}每天减少500$\sim$1000kcal能量摄入，每月减重1$\sim$2kg为宜。每日能量摄入不宜低于1000$\sim$1200kcal（女性）/ 1200$\sim$1600kcal（男性）。
    \item \textbf{调整宏量营养素供能比：}蛋白质\imp{20\%$\sim$25\%}，脂肪\imp{20\%$\sim$30\%}，碳水化合物\imp{45\%$\sim$60\%}。
    \item \textbf{增加膳食纤维：}每日\imp{25$\sim$30g}，增加饱腹感，延缓胃排空。
    \item \textbf{限制添加糖和饱和脂肪：}添加糖$<$总能量10\%，饱和脂肪酸$<$总能量10\%。
    \item \textbf{增加体力活动：}每周至少150分钟中等强度有氧运动 + 每周2$\sim$3次抗阻训练。推荐更高水平（每周200$\sim$300min）以维持体重下降和防止反弹。
    \item \textbf{行为矫正：}自我监测体重、记录饮食日记、规律进餐、三餐合理分配。
\end{enumerate}

\textbf{减重目标：}初始目标为减少原体重的5\%$\sim$10\%，即使未达到理想体重，减轻5\%已可显著改善血糖、血脂和血压等代谢指标。
\end{keybox}

% =====================================================================
%  体重控制与体力活动指南
% =====================================================================
\section{体重控制与体力活动指南}

\subsection{体重控制的意义及方法}

\begin{defbox}
\textbf{体重控制的意义（熟悉）：}
超重和肥胖是心血管疾病、2型糖尿病、多种癌症、骨关节炎、睡眠呼吸暂停等慢性病的重要危险因素。体重控制不仅改善外观，更重要的是降低慢性病风险、改善代谢指标、提高生活质量。即使减轻原体重的5\%-10\%，也可显著改善血糖、血脂和血压等代谢指标。

\textbf{体重控制的方法（熟悉）：}
\begin{enumerate}[leftmargin=*]
    \item \textbf{饮食控制：}控制总能量摄入——每日减少500-1000 kcal，每月减重1-2 kg为宜。女性不低于1000-1200 kcal/d，男性不低于1200-1600 kcal/d。
    \item \textbf{增加体力活动：}有氧运动结合抗阻训练（见体力活动指南）。
    \item \textbf{行为矫正：}自我监测（体重、饮食日记）、规律进餐（三餐合理分配）、控制进食速度、减少外出就餐。
    \item \textbf{综合干预：}饮食+运动+行为矫正联合干预效果优于单一措施。
    \item \textbf{药物和手术治疗：}BMI $\ge$ 28或BMI $\ge$ 24且伴有合并症者，在上述生活方式干预无效时可考虑。
\end{enumerate}
\end{defbox}

\subsection{体力活动指南}

\begin{defbox}
\textbf{体力活动指南（熟悉）：}

\textbf{一般健康人群的身体活动建议：}
\begin{enumerate}[leftmargin=*]
    \item \textbf{有氧运动：}每周至少150分钟中等强度有氧运动（如快走、骑车、游泳），或每周至少75分钟高强度有氧运动（如跑步）。每次持续至少10分钟。
    \item \textbf{抗阻训练：}每周至少2次（非连续日），涉及主要肌肉群（腿、臀、背、腹、胸、肩、臂）。
    \item \textbf{避免久坐：}任何强度的身体活动都有益，减少久坐时间。
\end{enumerate}

\textbf{控制体重的身体活动建议（更高水平）：}
\begin{itemize}[leftmargin=*]
    \item 建议每周增加有氧运动至\textbf{200-300分钟}中等强度。
    \item 维持体重下降和防止减重后反弹即需要达此水平。
\end{itemize}

\textbf{儿童青少年身体活动指南：}
\begin{itemize}[leftmargin=*]
    \item 每天至少\textbf{60分钟}中高强度身体活动。
    \item 每周至少3次高强度活动和力量训练。
    \item 视屏时间每天不超过2小时。
\end{itemize}
\end{defbox}

% =====================================================================
%  营养与痛风
% =====================================================================
\section{营养与痛风}

\begin{defbox}
\textbf{痛风（Gout）：}
痛风是由嘌呤代谢紊乱和/或尿酸排泄减少所致，以高尿酸血症为生化基础，以急性关节炎反复发作、痛风石形成、关节畸形、慢性间质性肾炎和尿酸性肾结石形成为临床特征的代谢性疾病。

\textbf{膳食因素：}
\begin{itemize}[leftmargin=*]
    \item 高嘌呤食物摄入过多（动物内脏、红肉、海鲜尤其贝类和沙丁鱼、肉汤）。
    \item 酒精（尤其啤酒和烈酒）——既增加尿酸生成又抑制尿酸排泄。
    \item 含糖饮料和果糖——促进尿酸生成。
    \item 水摄入不足——影响尿酸从肾脏排泄。
\end{itemize}
\end{defbox}

\begin{keybox}
\textbf{痛风的营养防治原则（掌握）：}

\begin{enumerate}[leftmargin=*]
    \item \textbf{限制总能量摄入：}超重/肥胖者需控制能量，减轻体重可降低血尿酸水平。但减重不宜过快（避免酮体升高抑制尿酸排泄）。
    \item \textbf{限制脂肪和蛋白质摄入：}脂肪供能比控制在20\%-25\%，过量脂肪可阻碍尿酸排泄。蛋白质不超过1.0 g/(kg$\cdot$d)，急性期0.8 g/(kg$\cdot$d)。
    \item \textbf{限制钠盐摄入：}低盐饮食（食盐$<$5 g/d）——有利于尿酸排泄。
    \item \textbf{\imp{限制嘌呤摄入：}}\\ 急性发作期：嘌呤摄入 < 150 mg/d（严格限制动物内脏、海鲜、肉汤、浓汤等）；\\ 缓解期：可适当放宽，但仍应控制高嘌呤食物。
    \item \textbf{增加蔬菜摄入：}多吃新鲜蔬菜水果，提供碱性矿物质和维生素C利于尿酸排泄。菠菜、花菜等虽含中等嘌呤但属植物性嘌呤，对血尿酸影响远小于动物性嘌呤。
    \item \textbf{保证充足饮水：}\imp{每日饮水2000-3000 ml}——稀释尿液，促进尿酸排泄，预防肾结石。
    \item \textbf{\imp{限制饮酒：}}尤其啤酒和烈酒，急性期应禁酒。
    \item \textbf{限制含果糖饮料：}果糖饮料可增加尿酸生成。
\end{enumerate}
\end{keybox}

% =====================================================================
%  第二节 营养与糖尿病
% =====================================================================
\section{营养与糖尿病}

\begin{defbox}
\textbf{糖尿病（Diabetes Mellitus）分类：}
\begin{itemize}[leftmargin=*]
    \item \textbf{1型糖尿病（T1DM）：}胰岛$\beta$细胞破坏，胰岛素绝对缺乏，多发生于儿童和青少年，需终身胰岛素治疗。
    \item \textbf{2型糖尿病（T2DM）：}胰岛素抵抗为主伴相对胰岛素缺乏，或胰岛素分泌缺陷为主伴胰岛素抵抗。占糖尿病患者90\%以上，与肥胖、不合理膳食和缺乏运动密切相关。
\end{itemize}
\end{defbox}

\begin{keybox}
\textbf{GI和GL在糖尿病饮食管理中的应用：}

\textbf{血糖生成指数（Glycemic Index, GI）：}反映食物摄入后升高血糖的速度和能力。
\begin{itemize}[leftmargin=*]
    \item \imp{低GI食物（GI $\leq$ 55）：}全谷物、豆类、牛奶
    \item \imp{中GI食物（GI 56$\sim$69）：}全麦面包、糙米
    \item \imp{高GI食物（GI $\geq$ 70）：}白面包、白米饭、含糖饮料
\end{itemize}
GI越低，血糖升高速度越慢。脂肪和蛋白质可延缓消化速度，因此脂肪和蛋白质含量高的食物GI较低。

\textbf{血糖负荷（Glycemic Load, GL）：}评价某种食物摄入量对人体血糖影响的幅度。\imp{GL $=$ GI $\times$ 食物中碳水化合物含量(g) / 100}。
\begin{itemize}[leftmargin=*]
    \item 低GL：GL $<$ 10；中GL：10$\sim$20；高GL：GL $>$ 20
    \item GI相同的食物，其GL可能不同。GL越高，餐后血糖总体水平越高。
\end{itemize}
\end{keybox}

\begin{keybox}
\textbf{糖尿病的饮食控制原则：}

\begin{enumerate}[leftmargin=*]
    \item \textbf{控制总能量：}以达到或维持理想体重为目标，根据BMI和活动水平个体化制定。
    \item \textbf{合理分配三大营养素：}
    \begin{itemize}[leftmargin=*]
        \item 碳水化合物：\imp{占总能量50\%$\sim$60\%}。优先选择低GI食物（全谷物、豆类），严格限制单糖和双糖。
        \item 蛋白质：\imp{占总能量15\%$\sim$20\%}。优质蛋白占1/3以上。有糖尿病肾病者需控制蛋白质摄入量。
        \item 脂肪：\imp{占总能量20\%$\sim$30\%}。饱和脂肪酸$<$7\%总能量，限制胆固醇$<$300mg/d。
    \end{itemize}
    \item \textbf{定时定量进餐：}三餐能量分配：早1/5、中2/5、晚2/5，或早中晚各1/3。少量多餐，避免血糖大幅波动。
    \item \textbf{增加膳食纤维：}\imp{每日25$\sim$30g}，可溶性纤维有助于降低餐后血糖。
    \item \textbf{综合管理：}饮食治疗是糖尿病\imp{最基本、最有效}的治疗措施，须与运动、药物治疗和血糖监测有机结合。
\end{enumerate}
\end{keybox}

% =====================================================================
%  第三节 营养与心血管疾病
% =====================================================================
\section{营养与心血管疾病}

\begin{keybox}
\textbf{膳食脂肪酸与血脂的关系：}

\begin{itemize}[leftmargin=*]
    \item \imp{饱和脂肪酸（SFA）：}HDL$\downarrow$，LDL$\uparrow$（LDL为心血管危险因素）→ 增加心血管疾病危险。主要来源：动物脂肪、棕榈油、椰子油。
    \item \imp{单不饱和脂肪酸（MUFA）：}HDL$\uparrow$，LDL$\downarrow$。代表：油酸（橄榄油、茶油）。对心血管有保护作用。
    \item \imp{多不饱和脂肪酸（PUFA）：}HDL$\downarrow$，LDL$\downarrow\downarrow$。代表：亚油酸（n-6）、$\alpha$-亚麻酸（n-3），存在于植物油中。虽强力降低LDL，但同时也降低HDL。
    \item \imp{反式脂肪酸：}HDL$\downarrow$，LDL$\uparrow$→ 增加心血管疾病危险。主要来源：氢化植物油、糕点、油炸食品。
\end{itemize}

\textbf{脂肪酸供能比例建议：}SFA:MUFA:PUFA $=$ 1:1:1。SFA $<$ 总能量7\%$\sim$10\%，反式脂肪酸摄入越少越好。
\end{keybox}

\begin{keybox}
\textbf{钠盐与高血压：}
\begin{itemize}[leftmargin=*]
    \item 高钠摄入是\imp{高血压最重要的膳食危险因素}。钠摄入增加→血容量增加→血压升高。
    \item \imp{每日食盐 $<$ 5g（约2000mg钠）}。减少加工食品、咸菜、调味品。
    \item \textbf{钾的降压作用：}钾可促进钠的排出，舒张血管，降低血压。多食蔬菜、水果、豆类和全谷物增加钾摄入。
    \item \textbf{钙、镁的作用：}钙和镁摄入充足也有助于血压控制。
    \item \textbf{膳食胆固醇：}\imp{$<$ 300mg/日}。减少动物内脏、脑、蛋黄等高胆固醇食物。
\end{itemize}
\end{keybox}

\begin{defbox}
\textbf{DASH饮食（Dietary Approaches to Stop Hypertension）：}高钾、高钙、高镁、高纤维、高蛋白，低饱和脂肪、低总脂肪、低胆固醇、低钠。强调蔬菜、水果、低脂奶制品、全谷物、鱼、禽肉和坚果，限制钠盐、红肉、甜食和含糖饮料。临床试验证实可\imp{降低收缩压8$\sim$14mmHg}。
\end{defbox}

\begin{keybox}
\textbf{动脉粥样硬化性心脏病的营养防治：}
\begin{enumerate}[leftmargin=*]
    \item 限制总能量摄入，维持理想体重
    \item 限制脂肪和胆固醇摄入，SFA $<$ 7\%$\sim$10\%总能量
    \item 提高植物性蛋白摄入，少吃甜食
    \item 摄入充足膳食纤维（25$\sim$30g/d）、矿物质和维生素
    \item 适当多吃富含植物化学物的食物
    \item 饮食清淡，少盐限酒
    \item 戒烟，保持积极生活方式
\end{enumerate}
\end{keybox}

% =====================================================================
%  第四节 营养与癌症
% =====================================================================
\section{营养与癌症}

\begin{keybox}
\textbf{膳食致癌因素：}

\begin{enumerate}[leftmargin=*]
    \item \textbf{高温烹调产生的致癌物：}
    \begin{itemize}[leftmargin=*]
        \item \imp{杂环胺（Heterocyclic Amines, HCAs）：}蛋白质含量较高的食物（肉类、鱼类）在高温烹调（烧、烤、煎、炸）时产生。加热温度愈高、时间愈长、水分含量愈少，产生的杂环胺愈多。美拉德反应在杂环胺形成中可能起催化作用。主要靶器官为肝脏，其次为血管、肠道、乳腺等。需经细胞色素P450代谢活化后才具有致突变性。
        \item \imp{多环芳烃（Polycyclic Aromatic Hydrocarbons, PAHs）：}代表物质为苯并(a)芘（B(a)P）。烘烤、烟熏食品中含量较高。B(a)P属于前致癌物，在体内经芳烃羟化酶作用代谢活化为多环芳烃环氧化物，与DNA、RNA和蛋白质等生物大分子结合而诱发肿瘤。
    \end{itemize}
    \item \textbf{腌制食品中的亚硝酸盐：}硝酸盐和亚硝酸盐用作食品防腐剂和护色剂，在体内可与胺类化合物反应生成\imp{亚硝胺（Nitrosamines）}——强致癌物，与食管癌、胃癌等有关。蔬菜从土壤中吸收硝酸盐也是体内硝酸盐的重要来源。
    \item \textbf{霉变食品中的黄曲霉毒素（Aflatoxin, AF）：}黄曲霉和寄生曲霉产生的代谢产物，\imp{黄曲霉毒素B$_1$为已知最强的化学致癌物之一}，主要诱发肝癌。最适产毒温度25$\sim$33$^\circ$C。花生最易受污染，其次为棉籽和油菜籽。
\end{enumerate}
\end{keybox}

\begin{keybox}
\textbf{膳食保护因素：}

\begin{enumerate}[leftmargin=*]
    \item \textbf{膳食纤维：}可吸附杂环胺并降低其活性，降低结直肠癌风险。增加粪便体积，稀释致癌物，缩短肠道转运时间。
    \item \textbf{抗氧化维生素：}维生素C、维生素E、$\beta$-胡萝卜素等可清除自由基，阻断亚硝胺体内合成，保护DNA免受氧化损伤。
    \item \textbf{植物化学物：}蔬菜水果中的酚类、黄酮类、有机硫化合物等成分有抑制杂环胺致突变性和致癌性的作用，可诱导解毒酶活性，抑制肿瘤细胞增殖。
\end{enumerate}
\end{keybox}

\begin{tipbox}
\textbf{世界癌症研究基金会（WCRF）核心建议：}保持健康体重；积极运动（每周至少150分钟）；摄入丰富的全谷物、蔬菜、水果和豆类；限制快餐和高脂/高淀粉/高糖加工食品；限制红肉（每周不超过350$\sim$500g）和加工肉类（加工肉类被IARC列为Ⅰ类致癌物）；限制含糖饮料；限制饮酒；不依赖膳食补充剂预防癌症。
\end{tipbox}

% =====================================================================
%  第五节 营养与骨质疏松
% =====================================================================
\section{营养与骨质疏松}

\begin{defbox}
\textbf{骨质疏松症（Osteoporosis）：}以骨量减少、骨组织微结构破坏、骨脆性增加和易于骨折为特征的代谢性骨病。\imp{绝经后骨质疏松症（PMOP）}是最常见的类型，与雌激素水平下降密切相关。
\end{defbox}

\begin{defbox}
\textbf{钙、维生素D、蛋白质对骨健康的影响：}

\textbf{1. 钙（Calcium）：}
\begin{itemize}[leftmargin=*]
    \item 钙是骨骼最主要的矿物成分，占人体钙总量的99\%。
    \item 长期钙摄入不足→骨钙动员→骨量丢失→骨质疏松。
    \item 膳食因素影响钙吸收：\imp{促进因素：}维生素D、乳糖、适量蛋白质。\imp{抑制因素：}草酸（菠菜）、植酸（谷物）、膳食纤维过多、高钠（钠$\uparrow$→尿钙$\uparrow$）、咖啡因。
    \item 奶制品是钙的最佳来源。
\end{itemize}

\textbf{2. 维生素D：}
\begin{itemize}[leftmargin=*]
    \item 促进肠道钙磷吸收，促进肾小管钙重吸收，调节骨钙动员和沉积。
    \item 缺乏→儿童佝偻病、成人骨软化症、老年人骨质疏松症。
    \item 来源：皮肤经紫外线照射合成（主要来源）、鱼肝油、蛋黄、强化食品。
\end{itemize}

\textbf{3. 蛋白质：}
\begin{itemize}[leftmargin=*]
    \item 适量的蛋白质是维持骨基质（胶原蛋白）合成所必需。
    \item \imp{蛋白质过剩（尤其动物蛋白过多）：}含硫氨基酸摄入过多→代谢产酸→骨钙动员缓冲→加速骨骼中钙的丢失→骨质疏松症风险增加。
\end{itemize}
\end{defbox}

\begin{defbox}
\textbf{峰值骨量（Peak Bone Mass）：}
\begin{itemize}[leftmargin=*]
    \item 人在生命早期（儿童青少年期至30岁左右）骨量达到的最高值。
    \item \imp{峰值骨量越高，晚年发生骨质疏松的风险越低。}提升峰值骨量是预防骨质疏松最重要的策略之一。
    \item 影响峰值骨量的主要因素：遗传（最主要）、钙摄入、维生素D水平、体力活动（尤其是负重运动）、性激素水平。
\end{itemize}

\textbf{骨质疏松的预防措施：}
\begin{enumerate}[leftmargin=*]
    \item 保证充足钙摄入：成人\imp{800$\sim$1000mg/d}，老年人/绝经后女性\imp{1000$\sim$1200mg/d}。
    \item 充足维生素D：多晒太阳，必要时补充维生素D制剂。
    \item 适量优质蛋白质，但避免过量动物蛋白摄入。
    \item 限制高钠饮食（减少尿钙排出）。
    \item 规律负重运动（步行、跑步、跳跃、力量训练），促进骨形成。
    \item 避免吸烟和过量饮酒。
    \item 儿童青少年期最大化峰值骨量，成年后维持骨量，老年期减缓骨丢失。
\end{enumerate}
\end{defbox}

% =====================================================================
%  临床营养
% =====================================================================
\section{临床营养筛查与营养评价}

\subsection{基本概念与营养风险筛查}

\begin{defbox}
\textbf{基本概念（熟悉）：}

\begin{enumerate}[leftmargin=*]
    \item \textbf{营养风险（Nutritional Risk）：}现存或潜在的与营养因素相关的导致病人出现不利临床结局的风险。
    \item \textbf{营养风险筛查（Nutritional Risk Screening）：}发现病人是否存在营养问题和是否需要进一步进行全面营养评估的过程。目的是发现个体是否存在营养不足和有营养不足的危险。
    \item \textbf{NRS 2002（Nutrition Risk Screening 2002）：}
    \begin{itemize}[leftmargin=*]
        \item 地位：目前以临床结局改善与否为目标的唯一风险筛查工具。
        \item 推荐：ESPEN、ASPEN、CSPEN等均给予推荐（回顾性循证研究、前瞻性队列研究验证）。
        \item 适用人群：18岁以上且住院时间超过24小时的患者，不推荐用于未成年人（2018年CSPEN专家共识）。
    \end{itemize}
    \item \textbf{营养评估（Nutritional Assessment）：}临床营养专业人员通过人体测量、生化检查、临床检查及综合营养评价方法等手段，对患者营养代谢和机体功能等进行检查和评定，以确定营养不良的类型及程度。
    \item \textbf{评估金标准：SGA（Subjective Global Assessment，主观全面营养评价）。}
\end{enumerate}

\textbf{膳食营养评价的基本方法：}
\begin{itemize}[leftmargin=*]
    \item 膳食调查（入院前/后摄食情况）、人体测量、生化检验、临床检查。
\end{itemize}

\textbf{营养评定的生化指标（了解）：}
\begin{itemize}[leftmargin=*]
    \item 维生素和微量元素：应激状态（手术、烧伤、败血症等）对维生素和微量元素需要量显著增加。
    \item 血浆氨基酸谱：重度PEM时血浆总氨基酸值明显下降；必需氨基酸（EAA）下降较非必需氨基酸（NEAA）更明显——Val、Leu、Ile和Met下降最多。
    \item 谷氨酰胺（Gln）是人体含量最丰富的氨基酸，精氨酸（Arg）可提高免疫活性、促进蛋白质合成。
\end{itemize}
\end{defbox}

% =====================================================================
\section{临床营养治疗与病人膳食}
% =====================================================================

\subsection{临床营养治疗目的和方法}

\begin{defbox}
\textbf{规范化营养支持治疗4步骤（熟悉）：}
\[
\text{营养筛查（Nutritional Screening）} \rightarrow \text{营养评定（Nutritional Assessment）} \rightarrow \text{营养干预（Nutritional Intervention）} \rightarrow \text{监测（Monitoring）}
\]

\textbf{营养治疗的五步治疗模式——临床营养干预阶梯（了解）：}
\begin{enumerate}[leftmargin=*]
    \item \textbf{首选}饮食营养治疗（膳食调整和营养教育）
    \item 进食不足 → 饮食营养治疗 + \textbf{口服营养补充（ONS）}
    \item 无法进食或只能进食流质 → \textbf{全肠内营养（TEN）}
    \item 消化功能或耐受性差 → \textbf{肠内营养 + 肠外营养（EN+PN）}
    \item 完全无法进食或无法通过胃肠提供营养 → \textbf{全肠外营养（TPN）}
\end{enumerate}
\end{defbox}

\subsection{病人膳食分类}

\begin{defbox}
\textbf{医院膳食分类（熟悉）：}

医院膳食可分为：\textbf{基本膳食（医院常规膳食）、治疗膳食、诊断试验膳食。}

\textbf{一、基本膳食（4种）：}

\begin{enumerate}[leftmargin=*]
    \item \textbf{普通膳食：}与健康人膳食基本相同，能量及营养素供应充足，符合平衡膳食原则。\\ 适用：体温正常或接近正常，无咀嚼或消化吸收功能障碍，无特殊膳食要求者。\\ 配膳原则：品种多样化，合理分配——早25\%-30\%，中40\%，晚30\%-35\%。

    \item \textbf{软食：}质地软、易咀嚼、少渣、易消化。\\ 适用：低热、消化不良或咀嚼困难者，老年患者或3-4岁患儿。

    \item \textbf{半流质膳食：}半流体、细软、更易咀嚼和消化，少食多餐。\\ 适用：发热较高、消化道疾病、耳鼻喉术后、身体虚弱者。\\ 配膳原则：能量低于软食，膳食纤维更少，少食多餐。

    \item \textbf{流质膳食：}极易消化，含渣很少，流体状态或在口腔内融化为液体。\\ 适用：高热、极度虚弱、无咀嚼能力、急性传染病、病情危重或术后，肠道手术术前准备。\\ 配膳原则：不平衡膳食（能量和营养素均不足），每日6-7餐。\\
    分类：清流质（食道/胃肠道大手术后）、浓流质（口腔手术后）、冷流质（扁桃体摘除后、胃肠道出血期）、不胀气流质（腹部手术后）。
\end{enumerate}
\end{defbox}

% =====================================================================
\section{肠内营养、肠外营养与FSMP}
% =====================================================================

\subsection{肠内营养（Enteral Nutrition, EN）}

\begin{defbox}
\textbf{肠内营养（熟悉）：}

\textbf{概念：}有胃肠道消化吸收功能的病人，因机体病理、生理改变或治疗特殊要求，需利用口服等方式给予要素膳制剂，经胃肠道消化吸收，提供能量和营养素，满足机体代谢需要的营养支持疗法。

\textbf{核心地位：}EN是最符合生理要求的途径，也是营养支持的首选途径。

\textbf{优势：}
\begin{itemize}[leftmargin=*]
    \item 营养素经胃肠消化吸收到肝脏，有利于内脏蛋白质合成和代谢调节。
    \item 实施简单、并发症少、促进胃肠功能恢复、促进胃肠激素释放。
    \item 改善门静脉循环、防止肠粘膜萎缩和细菌移位。
    \item 对技术和设备要求低，费用较低。
\end{itemize}

\textbf{局限性：}
\begin{itemize}[leftmargin=*]
    \item 受胃肠功能限制（肠梗阻、高位肠瘘、高流量肠瘘等）。
    \item 受消化功能限制（胃酸缺乏、胆汁分泌障碍、胰酶分泌障碍）。
    \item 受消化和吸收功能限制（吸收不良综合症、短肠综合症、先天性巨结肠等）。
\end{itemize}

\textbf{输入途径和方法：}
\begin{enumerate}[leftmargin=*]
    \item \textbf{经口营养 / ONS：}口服特殊制备的营养物质。当患者经口还能进食但进食量不足或结构不合理时采用。
    \item \textbf{管饲营养：}经鼻-胃管、鼻-十二指肠/空肠管、胃造瘘（PEG）、空肠造瘘等途径输入。
    \item \textbf{输注方式：}一次投给（每次250-400 ml，4-6次/天）、间歇重力滴注、经泵连续输注（12-24小时）。
\end{enumerate}

\textbf{适应证：}①无法经口摄食、摄食不足或有摄食禁忌者；②胃肠道疾病；③胃肠道外疾病。

\textbf{禁忌证：}肠道梗阻。

\textbf{肠内营养制剂（了解）：}
\begin{itemize}[leftmargin=*]
    \item 要素膳（单体膳）：氮源为氨基酸或短肽，无需消化可直接吸收，适用于消化功能严重受损者。
    \item 非要素膳（聚合膳）：氮源为整蛋白，需消化酶消化，适用于消化功能基本正常者。
    \item 组件膳：含单一或某类营养素，可根据病情需要调配。
    \item 特殊用途制剂：针对特定疾病状态（肝病、肾病、糖尿病、呼吸系统疾病等）。
\end{itemize}
\end{defbox}

\subsection{肠外营养（Parenteral Nutrition, PN）}

\begin{defbox}
\textbf{肠外营养（熟悉）：}

\textbf{概念：}通过肠道外的通路即静脉途径输注能量和各种营养素，达到纠正或预防营养不良、维持营养平衡目的的营养补充方式。

\textbf{输入途径：}
\begin{enumerate}[leftmargin=*]
    \item \textbf{周围静脉置管：}方便，可在普通病房实施。适用于短期肠外营养支持或作为肠内营养补充。局限性：需较多液体稀释以保证充足能量，不适合需限制液体的患者。（使用浅表静脉，多为上肢末梢静脉。）
    \item \textbf{中心静脉置管：}如输液港（PORT），尖端位于上腔静脉，适用长期肠外营养。
\end{enumerate}

\textbf{输注形式与方法：}
\begin{enumerate}[leftmargin=*]
    \item \textbf{持续输注：}24h内均匀输注，优点为血糖稳定，缺点为持续高胰岛素水平可能致脂肪肝和肝酶升高。
    \item \textbf{循环输注：}12-18h输注全天营养液，可预防肝毒性，改善生活质量。局限性：需心功能耐受短时间内大量液体。
\end{enumerate}

\textbf{适应证（了解）：}胃肠道功能障碍或衰竭的病人，存在营养不良或预计2周内无法正常饮食者。

\textbf{并发症（熟悉）：}
\begin{enumerate}[leftmargin=*]
    \item \textbf{置管并发症：}气胸、血胸、血肿、损伤胸导管/动脉/神经、空气栓塞等。
    \item \textbf{感染并发症：}导管性败血症是肠外营养最常见的严重并发症。
    \item \textbf{代谢并发症：}液体超负荷、糖代谢紊乱、肝功能损害、酸碱平衡失调、电解质紊乱、代谢性骨病。
    \item \textbf{肠道并发症：}肠道黏膜萎缩（长期不使用肠道）。
\end{enumerate}

\textbf{禁忌证（了解）：}严重循环/呼吸功能衰竭、水电解质平衡紊乱、肝肾衰竭等。

\textbf{肠外营养制剂（了解）：}
\begin{itemize}[leftmargin=*]
    \item 氨基酸注射液（氮源）、脂肪乳剂（能量+必需脂肪酸）、葡萄糖注射液（主要能源）、维生素制剂、微量元素制剂、电解质溶液。
    \item "全营养混合液（TNA）"即"全合一（All-in-One）"是将所有营养成分混合于一个容器中输注。
\end{itemize}
\end{defbox}

\subsection{特殊医学用途配方食品（FSMP）}

\begin{defbox}
\textbf{FSMP概念（熟悉）及分类（了解）：}

\textbf{概念：}特殊医学用途配方食品（Food for Special Medical Purposes, FSMP）是为满足进食受限、消化吸收障碍、代谢紊乱或特定疾病状态人群对营养素或膳食的特殊需要，专门加工配制而成的配方食品。必须在医生或临床营养师指导下，单独食用或与其他食品配合食用。

\textbf{FSMP不同于普通食品和保健食品：}FSMP不是药品，不能替代药物的治疗作用，但它是疾病治疗和康复过程中不可或缺的营养支持手段，属于肠内营养的范畴。

\textbf{分类（了解）：}
\begin{enumerate}[leftmargin=*]
    \item \textbf{全营养配方食品：}可作为单一营养来源满足目标人群营养需求。
    \item \textbf{特定全营养配方食品：}可作为单一营养来源满足目标人群在特定疾病或医学状态下的营养需求（如糖尿病全营养配方、肾病全营养配方、肿瘤全营养配方等）。
    \item \textbf{非全营养配方食品：}可满足目标人群部分营养需求，不适用于作为单一营养来源（如蛋白质组件、脂肪组件、碳水化合物组件、电解质配方、增稠组件等）。
\end{enumerate}

\textbf{管理：}我国对FSMP实行注册制管理，需经国家市场监督管理总局注册批准后方可生产销售。
\end{defbox}

% =====================================================================
%  复习题
% =====================================================================
\reviewcontent{
\section{复习题}
\begin{enumerate}[leftmargin=*, itemsep=8pt]
    \item \textbf{【名词解释】}BMI；中心性肥胖；GI；GL；DASH饮食；峰值骨量；痛风；NRS 2002；SGA；FSMP
    \item \textbf{【简答题】}简述膳食营养与慢性病（NCD）的关系。
    \item \textbf{【简答题】}简述体重控制的意义及方法。
    \item \textbf{【简答题】}简述体力活动指南对一般人群和体重维持的建议。
    \item \textbf{【简答题】}简述肥胖的判定标准和膳食营养防治策略。
    \item \textbf{【简答题】}简述GI和GL在糖尿病饮食管理中的应用。
    \item \textbf{【简答题】}简述糖尿病的饮食控制原则（至少五点）。
    \item \textbf{【简答题】}试述膳食脂肪酸（SFA、MUFA、PUFA、反式脂肪酸）与血脂的关系。
    \item \textbf{【简答题】}简述钠盐与高血压的关系及心血管疾病的膳食预防要点。
    \item \textbf{【简答题】}列举膳食致癌因素及其食物来源。
    \item \textbf{【简答题】}列举膳食保护因素及其防癌机制。
    \item \textbf{【简答题】}简述钙、维生素D、蛋白质对骨健康的影响。
    \item \textbf{【简答题】}简述峰值骨量的概念及骨质疏松的营养预防措施。
    \item \textbf{【简答题】}简述痛风的营养防治原则。
    \item \textbf{【简答题】}简述临床营养支持治疗的四步规范化流程和五步营养干预阶梯。
    \item \textbf{【简答题】}简述医院基本膳食的分类及各自适用对象。
    \item \textbf{【简答题】}简述肠内营养和肠外营养的概念、适应证、并发症，以及FSMP的概念和分类。
\end{enumerate}

\subsection*{参考答案要点}
\begin{enumerate}[leftmargin=*, itemsep=5pt]
    \item \textbf{BMI} = 体重(kg) / [身高(m)]$^2$，我国标准：BMI $\geq$ 28为肥胖，24.0$\sim$27.9为超重，18.5$\sim$23.9为正常。\textbf{中心性肥胖}：腰围男$\geq$85cm或女$\geq$80cm，或腰臀比男$>$0.9、女$>$0.85。\textbf{GI（血糖生成指数）：}反映食物摄入后升高血糖的速度和能力，低GI $\leq$55，中GI 56$\sim$69，高GI $\geq$70。\textbf{GL（血糖负荷）：}$=$ GI $\times$ 食物中碳水化合物含量(g) / 100，评价食物摄入量对血糖影响的幅度。\textbf{DASH饮食：}高钾/高钙/高镁/高纤维/高蛋白，低饱和脂肪/低总脂肪/低胆固醇/低钠的膳食模式，可降低收缩压8$\sim$14mmHg。\textbf{峰值骨量：}生命早期骨量达到的最高值（约30岁），峰值骨量越高，晚年骨质疏松风险越低。\textbf{痛风：}嘌呤代谢紊乱和/或尿酸排泄减少所致代谢病，以高尿酸血症为基础。\textbf{NRS 2002：}营养风险筛查2002，目前以临床结局改善为目标的唯一风险筛查工具。\textbf{SGA：}主观全面营养评价，是营养评估金标准。\textbf{FSMP：}特殊医学用途配方食品，为满足特定疾病人群特殊营养需要专门配制的配方食品。

    \item 膳食营养与NCD的关系：不合理膳食是NCD最重要可改变危险因素之一。高能量密度食物+含糖饮料等→肥胖→心血管病/糖尿病/多种癌症共同危险因素；高钠高饱和脂肪→心血管疾病；高GI/GL膳食→2型糖尿病；加工肉类（Ⅰ类致癌物）+酒精→癌症。合理膳食是NCD一级预防的基石。

    \item 体重控制意义：超重肥胖是心血管病、T2DM、多种癌症等的重要危险因素，减轻5\%-10\%体重即可显著改善代谢指标。方法：(1)饮食控制（日减500-1000kcal）；(2)增加体力活动（有氧+抗阻）；(3)行为矫正（自我监测、规律进餐）；(4)综合干预效果最佳；(5)严重肥胖者可考虑药物或手术。

    \item 体力活动指南——一般人群：每周$\ge$150min中等强度有氧或75min高强度+每周2次抗阻训练；体重维持：建议每周200-300min中等强度；儿童青少年：每天$\ge$60min中高强度活动+视屏时间$<$2h/d。

    \item \textbf{判定标准：}BMI $\geq$ 28为肥胖，24.0$\sim$27.9为超重；腰围男$\geq$85cm、女$\geq$80cm为中心性肥胖。\textbf{防治策略：}控制总能量（每日减少500$\sim$1000kcal），调整宏量营养素供能比（蛋白质20\%$\sim$25\%，脂肪20\%$\sim$30\%，碳水45\%$\sim$60\%），增加膳食纤维（25$\sim$30g/d），限制添加糖和饱和脂肪，配合有氧+抗阻运动（每周$\geq$150min中等强度），行为矫正（自我监测、规律进餐），初始减重目标为原体重的5\%$\sim$10\%，每月减重1$\sim$2kg。

    \item \textbf{GI的应用：}指导糖尿病患者选择低GI食物（全谷物、豆类、牛奶，GI $\leq$ 55），低GI有助于改善糖代谢，有效降低血糖水平；同时也可用于控制体重和慢性病。\textbf{GL的应用：}GL = GI $\times$ 碳水含量 / 100，GL $<$ 10为低GL、10$\sim$20为中GL、$>$ 20为高GL。GI相同的食物GL可能不同，GL越高餐后血糖总体水平越高，因此不仅要看GI，也要控制碳水化合物的摄入量。

    \item 五点原则：(1)控制总能量，以达到或维持理想体重为目标；(2)合理分配宏量营养素——碳水50\%$\sim$60\%（优先低GI食物）、蛋白质15\%$\sim$20\%（优质蛋白$>$1/3）、脂肪20\%$\sim$30\%（SFA $<$ 7\%）；(3)定时定量进餐，少量多餐，三餐能量分配早1/5、中2/5、晚2/5；(4)增加膳食纤维25$\sim$30g/d，可溶性纤维降低餐后血糖；(5)饮食治疗须与运动、药物治疗和血糖监测有机结合，饮食治疗是糖尿病最基本最有效的治疗。

    \item \textbf{SFA（饱和脂肪酸）：}HDL$\downarrow$、LDL$\uparrow$→增加心血管疾病危险，来源为动物脂肪、棕榈油。\textbf{MUFA（单不饱和脂肪酸）：}HDL$\uparrow$、LDL$\downarrow$→对心血管有保护作用，代表为油酸（橄榄油）。\textbf{PUFA（多不饱和脂肪酸）：}HDL$\downarrow$、LDL$\downarrow\downarrow$，强力降低LDL但也降低HDL，来源为植物油。\textbf{反式脂肪酸：}HDL$\downarrow$、LDL$\uparrow$→增加心血管疾病危险，来源为氢化植物油、糕点。建议SFA:MUFA:PUFA = 1:1:1。

    \item \textbf{钠盐与高血压：}高钠摄入是高血压最重要的膳食危险因素，钠$\uparrow$→血容量$\uparrow$→血压$\uparrow$，每日食盐$<$5g（约2000mg钠）。钾可促进钠排出、舒张血管而降低血压。\textbf{膳食预防要点：}控制能量维持理想体重，限制SFA（$<$7\%$\sim$10\%能量）和反式脂肪，限制胆固醇$<$300mg/d，控制钠摄入$<$5g/d盐，增加钾/钙/镁摄入，增加膳食纤维25$\sim$30g/d，限制饮酒，戒烟。

    \item \textbf{三类膳食致癌因素：}(1)\textbf{杂环胺：}高温烹调（烧/烤/煎/炸）蛋白质含量高的肉类产生，温度越高、时间越长、水分越少，产生越多，靶器官主要为肝脏。(2)\textbf{多环芳烃：}代表为苯并(a)芘，烟熏和烘烤食品中含量高，属前致癌物，经代谢活化后与DNA结合诱发肿瘤。(3)\textbf{亚硝酸盐：}腌制食品中的防腐剂和护色剂，与胺类反应生成亚硝胺（强致癌物），与食管癌、胃癌有关。(4)\textbf{黄曲霉毒素：}霉变食品，黄曲霉毒素B$_1$为最强化学致癌物之一，主要诱发肝癌，花生最易污染。

    \item \textbf{三大膳食保护因素：}(1)\textbf{膳食纤维：}吸附杂环胺降低其活性，稀释致癌物，缩短肠道转运时间，降低结直肠癌风险。(2)\textbf{抗氧化维生素（VC、VE、$\beta$-胡萝卜素）：}清除自由基，阻断亚硝胺体内合成，保护DNA。(3)\textbf{植物化学物（酚类、黄酮类等）：}抑制杂环胺致突变性和致癌性，诱导解毒酶活性，抑制肿瘤细胞增殖。

    \item \textbf{钙：}骨骼最主要矿物成分（占99\%），摄入不足→骨钙动员→骨质疏松。促进吸收因素：维生素D、乳糖、适量蛋白质；抑制因素：草酸、植酸、高钠、咖啡因。\textbf{维生素D：}促进肠道钙磷吸收和肾小管钙重吸收，调节骨代谢，缺乏→佝偻病/骨软化症/骨质疏松症。\textbf{蛋白质：}适量蛋白质维持骨基质合成所必需；蛋白质过剩（尤其动物蛋白过多）→含硫氨基酸产酸→骨钙动员缓冲→加速骨钙丢失。

    \item \textbf{峰值骨量：}生命早期（至30岁左右）骨量达到的最高值。峰值骨量越高，晚年骨质疏松风险越低。遗传为最主要影响因素。\textbf{预防措施：}保证充足钙摄入（成人800$\sim$1000mg/d，老年人1000$\sim$1200mg/d），充足维生素D，适量优质蛋白但避免过量，限制高钠饮食，规律负重运动，避免吸烟和过量饮酒。分三期策略：儿童青少年期最大化峰值骨量，成年后维持骨量，老年期减缓骨丢失。

    \item 痛风营养防治原则8条：(1)限制总能量摄入——减轻体重降低血尿酸，减重不宜过快；(2)限制脂肪——供能比20\%-25\%，过多阻碍尿酸排泄；(3)限制钠盐——$<$5g/d；(4)\textbf{限制嘌呤摄摄入——核心}：急性期$<$150mg/d（禁内脏海鲜肉汤），缓解期可放宽；(5)增加蔬菜摄入——提供碱性矿物质和VC；(6)\textbf{保证充足饮水——2000-3000ml/d}促尿酸排泄；(7)限酒——尤其啤酒烈酒，急性期禁酒；(8)限制含果糖饮料。

    \item 规范化营养支持治疗4步骤：营养筛查→营养评定→营养干预→监测。五步干预阶梯：(1)饮食治疗（首选）；(2)饮食+口服营养补充（ONS）；(3)全肠内营养（TEN）；(4)肠内+肠外（EN+PN）；(5)全肠外营养（TPN）。

    \item 医院基本膳食4类：普通膳食（体温正常无特殊要求）、软食（低热/咀嚼困难/老年患者）、半流质（发热较高/消化道疾病/术后）、流质（高热/极度虚弱/危重/术前，分为清流质、浓流质、冷流质、不胀气流质）。

    \item EN概念：通过胃肠途径给予要素膳制剂的营养支持方法，是营养支持首选途径（最符合生理）。适应证：无法经口摄食/摄食不足、胃肠道疾病、胃肠道外疾病。禁忌证：肠道梗阻。并发症：腹泻、腹胀、误吸。PN概念：通过静脉途径输注能量和营养素的营养补充方式。适应证：胃肠功能障碍或衰竭，存在营养不良或预计2周内无法正常饮食。并发症：置管并发症（气胸血胸）、感染（导管性败血症最常见严重并发症）、代谢并发症、肠道黏膜萎缩。FSMP：特殊医学用途配方食品，为满足特定疾病人群特殊营养需要专门配制的配方食品，须在医生/临床营养师指导下使用。分类：全营养配方、特定全营养配方、非全营养配方。
\end{enumerate}

}
\newpage
